\chapter{解析函数}

\problem{5.1}{请说明下面复变函数在复平面上的哪些点或区域上可导，在复平面上的哪些区域上解析。
\begin{enumerate}
  \item $f(z)=\operatorname{Re} z$ \qquad
  \item $f(z)=z^*$ \qquad
  \item $f(z)=|z|$
  \item $f(z)=\mathrm{e}^z$ \qquad
  \item $f(z)=z^2$
  \item $f(z)=(x^2-y^2-x)+\mathrm{i}(2xy-y^2)$ \qquad
  \item $f(z)=x^3-y^3-\mathrm{i}2x^2y$
\end{enumerate}}
\answerarea

\problem{5.2}{若 $f(z)=u(x,y)+\mathrm{i}v(x,y)$ 和 $f^*(z)=u(x,y)-\mathrm{i}v(x,y)$ 都为解析函数，请说明 $f(z)$ 的函数形式。}
\answerarea
\clearpage

\problem{5.3}{若 $f(z)=u(x,y)+\mathrm{i}v(x,y)$ 为解析函数，请证明 $f^*(z^*)$ 也是解析函数。}
\answerarea

\problem{5.4}{请证明：柱坐标系下的解析函数
\[
  f(z) = u(\rho,\varphi) + \mathrm{i}v(\rho,\varphi)
\]
满足的 Cauchy-Riemann 方程
\[
  \left\{
  \begin{aligned}
    \frac{\partial u}{\partial \rho} &= \frac{1}{\rho}\frac{\partial v}{\partial \varphi} \\
    \frac{1}{\rho}\frac{\partial u}{\partial \varphi} &= -\frac{\partial v}{\partial \rho}
  \end{aligned}
  \right.
\]}
\answerarea
\clearpage

\problem{5.5}{已知解析函数 $f(z)=u(x,y)+\mathrm{i}v(x,y)$ 的实部 $u(x,y)$ 或虚部 $v(x,y)$，求此解析函数。
\begin{enumerate}
  \item $v(x,y)=xy$ \qquad
  \item $v(x,y)=\dfrac{y}{x^2+y^2}$
  \item $u(x,y)=\mathrm{e}^{-y}\sin x$ \qquad
  \item $u(x,y)=\cos x\cosh y$
\end{enumerate}}
\answerarea

\problem{5.6}{在平面极坐标系（$z=\rho\mathrm{e}^{\mathrm{i}\varphi}$）中，已知解析函数 $f(z)=u(\rho,\varphi)+\mathrm{i}v(\rho,\varphi)$ 的实部 $u(\rho,\varphi)=\varphi$，且 $f(1)=0$，求此解析函数。}
\answerarea
\clearpage

\problem{5.7}{在直角坐标系中，已知解析函数 $f(z)=u(x,y)+\mathrm{i}v(x,y)$ 的实部和虚部之和为 $u(x,y)+v(x,y)=-(x+y)(x^2-4xy+y^2)$，求此解析函数。}
\answerarea

\problem{5.8}{有两条平行相距为 $2a$ 且均匀带电的细金属线，每单位长度所带的电荷量分别是 $+q$ 和 $-q$，求与此两金属线垂直的平面静电场的复势、电场线和等势线。}
\answerarea
\clearpage

\sproblem{5.9}{试分析以下几种常见的解析函数变换。
\begin{enumerate}
  \item 线性变换：$W=f(z)=az+b,\quad a,b\in R,\text{ 且 }a\neq 0$；
  \item 幂函数变换：$W=f(z)=z^n,\quad n\in\mathbf{Z}$；
  \item 根式变换：$W=f(z)=z^{\frac{1}{n}},\quad n\in\mathbf{Z}$；
  \item 对数变换：$W=f(z)=\ln z$；
  \item 指数变换：$W=f(z)=\mathrm{e}^z$；
  \item 根式变换：$W=f(z)=\dfrac{az+b}{cz+d},\quad a,b,c,d\in R\text{ 且 }ad-bc\neq 0$。
\end{enumerate}}
\answerarea
\clearpage
